\chapter{Educación}

\marginal{la función no se mueve}
La función Educación del Gobierno Federal, en bruto, pasa de \textbf{1,032,621.4 a
1,084,590.9 millones de pesos}: un aumento real de \textbf{0.7\%} y exactamente
\textbf{3.00\% del PIB} en los dos años. En un presupuesto donde el gasto
programable cae 4.1\% real, mantenerse plano es una decisión.

Es también el rubro más federalizado del paquete y por eso la primera pregunta no
es cuánto, sino quién lo ejerce.

\section{Federal contra federalizado}

\begin{table}[htbp]
\centering
\caption{Función Educación por ramo y subfunción (mdp)}
\label{tab:C7_1}
\small
\begin{tabular}{lSSSSS}
\toprule
{Concepto} & {2024} & {2025} & {Dif. nominal} & {Dif. real} & {Var. real \%} \\
\midrule
{TOTAL funcion Educacion (GF, bruto)} & \num{1032621.4} & \num{1084590.9} & \num{51969.4} & \num{7566.7} & \num{0.7} \\
\multicolumn{6}{l}{\itshape por ramo} \\
{33 Aportaciones Federales para Entidades Federativas y Municipios} & \num{526250.1} & \num{553261.9} & \num{27011.8} & \num{4383.1} & \num{0.8} \\
{11 Educación Pública} & \num{411848.3} & \num{437675.4} & \num{25827.1} & \num{8117.6} & \num{1.9} \\
{25 Previsiones y Aportaciones para los Sistemas de Educación Básica, Normal, Tecnológica y de Adultos} & \num{82460.5} & \num{81610.1} & \num{-850.5} & \num{-4396.3} & \num{-5.1} \\
{08 Agricultura y Desarrollo Rural} & \num{5346.7} & \num{5534.2} & \num{187.5} & \num{-42.4} & \num{-0.8} \\
{07 Defensa Nacional} & \num{3854.1} & \num{3992.7} & \num{138.5} & \num{-27.2} & \num{-0.7} \\
{13 Marina} & \num{2230.0} & \num{2138.2} & \num{-91.8} & \num{-187.7} & \num{-8.1} \\
{47 Entidades no Sectorizadas} & \num{631.7} & \num{378.4} & \num{-253.3} & \num{-280.4} & \num{-42.6} \\
\multicolumn{6}{l}{\itshape por subfuncion} \\
{subfuncion 1} & \num{650707.5} & \num{690468.9} & \num{39761.4} & \num{11781.0} & \num{1.7} \\
{subfuncion 3} & \num{165466.5} & \num{171539.1} & \num{6072.7} & \num{-1042.4} & \num{-0.6} \\
{subfuncion 2} & \num{145417.1} & \num{151385.4} & \num{5968.3} & \num{-284.6} & \num{-0.2} \\
{subfuncion 6} & \num{55256.3} & \num{55035.7} & \num{-220.7} & \num{-2596.7} & \num{-4.5} \\
{subfuncion 4} & \num{10322.0} & \num{10521.3} & \num{199.3} & \num{-244.5} & \num{-2.3} \\
{subfuncion 5} & \num{5452.1} & \num{5640.5} & \num{188.4} & \num{-46.1} & \num{-0.8} \\
\bottomrule
\end{tabular}
\par\vspace{2pt}\footnotesize Fuente: elaboración propia del ITED con los analíticos del PEF aprobado 2024 y 2025. Las diferencias en millones de pesos se reportan dos veces, nominal y real; la real está en pesos de 2025 con el deflactor de 4.3\% que declara el CGPE.
\end{table}


De los 1,084,590.9 millones, \textbf{el Ramo 33 pone 553,261.9, el 51\%}, y el
Ramo 11 pone 437,675.4, el 40\%. El Ramo 25, que financia la educación básica de
la Ciudad de México y las previsiones salariales, pone 81,610.1. Los ramos
militares, Agricultura y Entidades no Sectorizadas ponen el resto.

\textbf{Más de la mitad del gasto educativo federal no lo ejerce la Secretaría de
Educación Pública: se transfiere a las entidades federativas.} Es la distinción
que decide la lectura del rubro y ninguna pieza del paquete la presenta.

Por subfunción, la educación básica sube 1.7\% real y todo lo demás cae: media
superior 0.2\%, superior 0.6\%, posgrado 2.3\%, educación para adultos 0.8\% y
otros servicios educativos 4.5\%. \textbf{El único nivel que gana en términos
reales es el básico}, y lo hace por dos vías que el resto del capítulo separa.

\section{El FONE}

El Fondo de Aportaciones para la Nómina Educativa y el Gasto Operativo pasa de
496,792.7 a \textbf{521,406.5 millones de pesos}, un aumento real de 0.6\%. Es
\textbf{el 48\% de toda la función educativa} y equivale a 1.12 veces el Ramo 11
completo.

\footnotesize
\setlength{\tabcolsep}{4pt}
\begin{longtable}{L{4.2cm}SSSSS}
\caption{FONE por entidad federativa (mdp)}\label{tab:C7_3}\\
\toprule
{Concepto} & {2024} & {2025} & {Dif. nominal} & {Dif. real} & {Var. real \%} \\
\midrule\endfirsthead
\toprule
{Concepto} & {2024} & {2025} & {Dif. nominal} & {Dif. real} & {Var. real \%} \\
\midrule\endhead
\bottomrule\endfoot
{FONE TOTAL} & \num{496792.7} & \num{521406.5} & \num{24613.8} & \num{3251.7} & \num{0.6} \\
\multicolumn{6}{l}{\itshape por entidad federativa} \\
{15 Estado de México} & \num{49009.1} & \num{46705.0} & \num{-2304.1} & \num{-4411.5} & \num{-8.6} \\
{30 Veracruz} & \num{36335.3} & \num{35935.7} & \num{-399.5} & \num{-1962.0} & \num{-5.2} \\
{20 Oaxaca} & \num{31126.2} & \num{30853.4} & \num{-272.8} & \num{-1611.3} & \num{-5.0} \\
{14 Jalisco} & \num{24740.3} & \num{26629.8} & \num{1889.6} & \num{825.7} & \num{3.2} \\
{12 Guerrero} & \num{24412.3} & \num{26504.6} & \num{2092.3} & \num{1042.6} & \num{4.1} \\
{07 Chiapas} & \num{24657.5} & \num{26183.8} & \num{1526.3} & \num{466.1} & \num{1.8} \\
{16 Michoacán} & \num{23857.4} & \num{25406.0} & \num{1548.5} & \num{522.7} & \num{2.1} \\
{21 Puebla} & \num{22799.5} & \num{24339.4} & \num{1539.9} & \num{559.5} & \num{2.4} \\
{11 Guanajuato} & \num{20543.9} & \num{21909.6} & \num{1365.7} & \num{482.3} & \num{2.3} \\
{13 Hidalgo} & \num{18941.7} & \num{20612.9} & \num{1671.2} & \num{856.7} & \num{4.3} \\
{19 Nuevo León} & \num{18703.0} & \num{20045.2} & \num{1342.3} & \num{538.0} & \num{2.8} \\
{28 Tamaulipas} & \num{16634.5} & \num{17839.5} & \num{1205.1} & \num{489.8} & \num{2.8} \\
{08 Chihuahua} & \num{16202.3} & \num{17459.3} & \num{1257.1} & \num{560.4} & \num{3.3} \\
{24 San Luis Potosí} & \num{15114.5} & \num{16415.1} & \num{1300.6} & \num{650.6} & \num{4.1} \\
{02 Baja California} & \num{14267.3} & \num{15115.1} & \num{847.8} & \num{234.3} & \num{1.6} \\
{25 Sinaloa} & \num{13856.9} & \num{14780.4} & \num{923.6} & \num{327.7} & \num{2.3} \\
{05 Coahuila} & \num{13603.9} & \num{14526.9} & \num{923.0} & \num{338.1} & \num{2.4} \\
{26 Sonora} & \num{11304.9} & \num{12080.1} & \num{775.2} & \num{289.1} & \num{2.5} \\
{10 Durango} & \num{10142.5} & \num{11013.9} & \num{871.4} & \num{435.3} & \num{4.1} \\
{17 Morelos} & \num{9535.4} & \num{10279.9} & \num{744.5} & \num{334.5} & \num{3.4} \\
{27 Tabasco} & \num{9670.4} & \num{10267.9} & \num{597.5} & \num{181.7} & \num{1.8} \\
{32 Zacatecas} & \num{9187.4} & \num{9876.1} & \num{688.7} & \num{293.7} & \num{3.1} \\
{22 Querétaro} & \num{8814.3} & \num{9450.5} & \num{636.2} & \num{257.1} & \num{2.8} \\
{31 Yucatán} & \num{8202.8} & \num{8806.4} & \num{603.7} & \num{250.9} & \num{2.9} \\
{01 Aguascalientes} & \num{7798.6} & \num{8355.9} & \num{557.2} & \num{221.9} & \num{2.7} \\
{23 Quintana Roo} & \num{7362.5} & \num{7849.6} & \num{487.1} & \num{170.5} & \num{2.2} \\
{29 Tlaxcala} & \num{7133.5} & \num{7650.1} & \num{516.6} & \num{209.9} & \num{2.8} \\
{18 Nayarit} & \num{6835.0} & \num{7386.7} & \num{551.7} & \num{257.8} & \num{3.6} \\
{03 Baja California Sur} & \num{5930.7} & \num{6293.9} & \num{363.1} & \num{108.1} & \num{1.7} \\
{04 Campeche} & \num{5613.6} & \num{6028.5} & \num{414.9} & \num{173.5} & \num{3.0} \\
{06 Colima} & \num{4413.2} & \num{4761.2} & \num{348.0} & \num{158.2} & \num{3.4} \\
{34 No Distribuible Geográficamente} & \num{42.6} & \num{44.2} & \num{1.6} & \num{-0.2} & \num{-0.4} \\
\end{longtable}
\par\vspace{2pt}\footnotesize Fuente: elaboración propia del ITED con los analíticos del PEF aprobado 2024 y 2025. Las diferencias en millones de pesos se reportan dos veces, nominal y real; la real está en pesos de 2025 con el deflactor de 4.3 \% que declara el CGPE.
\normalsize


El reparto por entidad federativa es el corte que un capítulo de educación
necesita y que el género no publica. \textbf{La fórmula de distribución está en el
artículo 27 de la Ley de Coordinación Fiscal y no en el paquete}, de modo que este
documento puede mostrar el resultado del reparto pero no puede decir si la fórmula
cambió. Se declara la limitación.

Una nota de perímetro que se repite todos los años y conviene fijar: \textbf{la
Ciudad de México no recibe FONE}. Su educación básica y normal se financia por el
Ramo 25, con dos programas propios. Cualquier comparación entre entidades que no
lo advierta encontrará un cero donde hay un sistema educativo completo.

\section{Las becas}

Los cuatro programas de beca del Ramo 11 pasan de 102,958.0 a \textbf{133,615.1
millones de pesos}, y su peso dentro del ramo sube de 23.5 a \textbf{28.7\%}. Es
el movimiento de política educativa más claro del paquete.

Casi todo se concentra en un programa: \textbf{la beca universal de educación
básica pasa de 49,869.8 a 78,840.7 millones}, un aumento real de 26,826.5. En 2025
lleva el nombre de Rita Cetina. Contra ese aumento, tres programas emblemáticos
retroceden en términos reales: los servicios de educación media superior pierden
5,353.4 millones de pesos de 2025, el subsidio a universidades públicas estatales
4,754.0 y La Escuela es Nuestra 4,577.7, esta última con una caída nominal de
28,358.3 a 25,000.0.

\marginal{una beca contra tres programas}
La lectura es directa y no requiere adjetivos: \textbf{el paquete 2025 reasigna
dentro de la educación básica hacia una transferencia monetaria a las familias, y
lo financia con los servicios de media superior, el subsidio universitario y la
infraestructura escolar.} El programa de infraestructura social del sector
educativo desaparece: pasa de 1,170.0 a cero.

\section{Comparación estructural}

La razón adoptada para este capítulo es \emph{función Educación sobre pensiones y
jubilaciones}, que compara las dos grandes transferencias públicas de los extremos
opuestos del ciclo de vida con perímetros declarados.

\begin{itemize}
\item Contra pensiones contributivas: 0.709 en 2023, 0.689 en 2024, \textbf{0.662
en 2025}.
\item Contra pensiones totales, con las no contributivas: 0.558, 0.518,
\textbf{0.501}.
\end{itemize}

\textbf{Por cada peso que el paquete destina a pensiones contributivas, destina 66
centavos a educación; hace dos años eran 71.} Con las pensiones no contributivas
dentro, la relación cruza el umbral simbólico: \textbf{el gasto educativo es ahora
exactamente la mitad del gasto pensionario.} La serie cae en las dos versiones y
en los tres años, y no por un recorte educativo, que no lo hubo, sino porque
pensiones crece 4.7\% real mientras educación crece 0.7\%.

\section{El supuesto que no está}

\textbf{Ninguna pieza del paquete declara matrícula, cobertura, planteles ni
docentes.} No el CGPE, no la iniciativa de ingresos, no el decreto. En
consecuencia este documento \textbf{no publica ninguna cifra de gasto por
alumno}, y no por descuido: dividir 1,084,590.9 millones de pesos entre una
matrícula que nadie declaró produciría un número con la apariencia de un dato.

La ausencia no es neutral. Educación es el único rubro grande cuyo denominador
demográfico se contrae, de modo que un gasto plano en términos reales
\textbf{sube} por alumno si la matrícula cae y baja si crece. El paquete no
permite decir cuál. El capítulo 13 vuelve sobre esto porque es, a nuestro juicio,
el hueco más importante del Paquete Económico 2025.

\clearpage
